% WORKING NOTES --- the living scratch document. Anything goes: half-results, ideas we are chasing,
% things we tried, open questions, stuff we are unsure will reach the paper. Newest entry on top.
% Graduation flow: working_notes -> progress_*.tex (checkpoints) -> sections/ (final paper).
% Shares the bib + numbers macros, so citations and figures are consistent with the paper.
\documentclass{article}
% Shared preamble for ALL root documents in this project (paper, working notes, lit review,
% pipeline spec, progress reports). Each root does: \documentclass{article} % Shared preamble for ALL root documents in this project (paper, working notes, lit review,
% pipeline spec, progress reports). Each root does: \documentclass{article} % Shared preamble for ALL root documents in this project (paper, working notes, lit review,
% pipeline spec, progress reports). Each root does: \documentclass{article} \input{shared/macros}
% Keeps packages, notation, and the single-source-of-truth numbers consistent everywhere.
%
% NOTE: the final paper will later swap \documentclass{article} for the venue template (acmart for
% FAccT). This preamble stays largely reusable; drop the package lines the template already provides.

\usepackage[utf8]{inputenc}
\usepackage[T1]{fontenc}
\usepackage[margin=1in]{geometry}
\usepackage{amsmath,amssymb}
\usepackage{graphicx}
\usepackage{booktabs}
\usepackage[numbers]{natbib}  % numeric mode (plainnat); final paper swaps to the venue template's citations
\usepackage{xcolor}
\usepackage{hyperref}

% --- notation / shorthands ---------------------------------------------------------------------
\newcommand{\rewardmodel}{reward model}
\newcommand{\RMs}{reward models}
\newcommand{\diffmean}{difference-of-means}
\newcommand{\nullspace}{null-space projection}
\newcommand{\autoinf}{auto-influence}
\newcommand{\crossinf}{cross-influence}
\newcommand{\leace}{\textsc{Leace}}
\newcommand{\taco}{\textsc{TaCo}}
\newcommand{\todo}[1]{\textcolor{red}{[\textbf{TODO}: #1]}}

% --- single source of truth for numbers (see export_paper_numbers.py in the code repo) ----------
\input{shared/numbers}         % AUTO-GENERATED from results JSONs
\input{shared/numbers_manual}  % hand-maintained / not-yet-serialized

% Keeps packages, notation, and the single-source-of-truth numbers consistent everywhere.
%
% NOTE: the final paper will later swap \documentclass{article} for the venue template (acmart for
% FAccT). This preamble stays largely reusable; drop the package lines the template already provides.

\usepackage[utf8]{inputenc}
\usepackage[T1]{fontenc}
\usepackage[margin=1in]{geometry}
\usepackage{amsmath,amssymb}
\usepackage{graphicx}
\usepackage{booktabs}
\usepackage[numbers]{natbib}  % numeric mode (plainnat); final paper swaps to the venue template's citations
\usepackage{xcolor}
\usepackage{hyperref}

% --- notation / shorthands ---------------------------------------------------------------------
\newcommand{\rewardmodel}{reward model}
\newcommand{\RMs}{reward models}
\newcommand{\diffmean}{difference-of-means}
\newcommand{\nullspace}{null-space projection}
\newcommand{\autoinf}{auto-influence}
\newcommand{\crossinf}{cross-influence}
\newcommand{\leace}{\textsc{Leace}}
\newcommand{\taco}{\textsc{TaCo}}
\newcommand{\todo}[1]{\textcolor{red}{[\textbf{TODO}: #1]}}

% --- single source of truth for numbers (see export_paper_numbers.py in the code repo) ----------
% AUTO-GENERATED by experiments/export_paper_numbers.py -- DO NOT EDIT BY HAND.
% source: /Users/CedricESMT/AI_safety/AI_Safety_Frankfurt/greenTeam/OneBiasAfterAnotherFork/artifacts/results/demographic/*.json | generated: 2026-07-02 | code commit: 63be6dc
% Hand-maintained / not-yet-serialized numbers live in shared/numbers_manual.tex.

\newcommand{\autoInflIntersectionCVbase}{1.00}
\newcommand{\autoInflIntersectionCVnull}{0.01}
\newcommand{\autoInflIntersectionCreditbase}{1.00}
\newcommand{\autoInflIntersectionCreditnull}{0.00}
\newcommand{\autoInflSexCVbase}{1.00}
\newcommand{\autoInflSexCVnull}{0.11}
\newcommand{\autoInflSexCreditbase}{1.00}
\newcommand{\autoInflSexCreditnull}{0.06}
\newcommand{\batteryAgeExplicitCV}{0.23}
\newcommand{\batteryAgeProxyCV}{0.99}
\newcommand{\batteryFamilyExplicitCV}{0.04}
\newcommand{\batteryFamilyProxyCV}{0.91}
\newcommand{\batterySexProxyCV}{0.63}
\newcommand{\cosCorrectnessPLvsCommute}{0.94}
\newcommand{\crossAccBaselineCV}{0.59}
\newcommand{\crossAccBaselineCredit}{0.46}
\newcommand{\decisionDiscWinAge}{0.00}
\newcommand{\decisionDiscWinFamily}{0.39}
\newcommand{\decisionDiscWinIntersection}{0.32}
\newcommand{\decisionDiscWinSex}{0.00}
\newcommand{\decisionGapFairDiscAge}{+2.51}
\newcommand{\decisionGapFairDiscFamily}{+0.27}
\newcommand{\decisionGapFairDiscIntersection}{+0.10}
\newcommand{\decisionGapFairDiscSex}{+4.92}
\newcommand{\erasureLinearLeaceCorrPL}{52\%}
\newcommand{\erasureLinearNoneCorrPL}{100\%}
\newcommand{\erasureMlpLeaceCorrPL}{100\%}
\newcommand{\meanGapIntersectionCV}{+0.62}
\newcommand{\meanGapIntersectionCredit}{+0.75}
\newcommand{\meanGapSexCV}{+0.59}
\newcommand{\meanGapSexCredit}{+0.39}
\newcommand{\modelSmall}{Skywork/Skywork-Reward-V2-Qwen3-0.6B}
\newcommand{\nEval}{200}
\newcommand{\probeConclAccHeldoutCommute}{100\%}
\newcommand{\probeConclAccHeldoutPL}{100\%}
\newcommand{\probeCorrAccHeldoutCommute}{100\%}
\newcommand{\probeCorrAccHeldoutPL}{100\%}
\newcommand{\reasonConclusionCommute}{+1.46}
\newcommand{\reasonConclusionIntersection}{+0.92}
\newcommand{\reasonConclusionPL}{+0.67}
\newcommand{\reasonCorrectnessCommute}{+1.46}
\newcommand{\reasonCorrectnessIntersection}{+1.33}
\newcommand{\reasonCorrectnessPL}{+1.45}
\newcommand{\transferCommuteToPLbase}{+1.45}
\newcommand{\transferCommuteToPLnull}{+0.17}
         % AUTO-GENERATED from results JSONs
% HAND-MAINTAINED numbers: values not yet serialized to JSON, or human-derived.
% Everything auto-generated lives in shared/numbers.tex (do not edit that one by hand).
%
% Additivity cosines: additivity_check.py now supports --out; until it is re-run with --out these
% stay here. Replace with the auto macros (\additivityCosCV / \additivityCosCredit) once serialized.
\newcommand{\additivityCosCV}{0.66}     % TODO auto: run_additivity.py --domain cv --out ...
\newcommand{\additivityCosCredit}{0.71} % TODO auto: run_additivity.py --domain credit --out ... (STALE, see below)

% --- STALE AUTO-NUMBERS (2026-09-16) ---------------------------------------------------------------
% numbers.tex is auto-generated and must not be edited, so stale values are flagged here instead: each
% \markstale{name}{mark} re-typesets the macro in orange with a mark. Every use in every document is
% flagged at once. DELETE THE MATCHING LINE once export_paper_numbers.py has re-exported the value from
% a re-run (a leftover line would flag a fresh number as stale).
%   \dag      credit: German Credit was decoded with the wrong UCI code table, and the pairs now come
%             from the sex x age x marital-status factorial -- every credit result must be re-run.
%   \ddag     ASAP: redaction-tag cleaning changed 22 of 12,976 essays (adjacent tags such as
%             @PERCENT1@NUM1) -- expected to be negligible, re-run for consistency.
%   \S        inputs changed: the proxy first names are now Haim et al.'s Table 3 (every name-proxy
%             number), and the decision-response verdicts now open without the synthetic "N years" field.
%   \P        education substrate changed: (a) the essay-id lookup keyed on the lossy `essay_id` column
%             and merged two essays, it now keys on `essay_id_comp`, so every PERSUADE record id and the
%             sampled set change; (b) the essays are now restricted to the stage-neutral prompts without
%             in-essay pupil cues; (c) the grade_level poles widened to 6th grade vs doctoral candidate.
%             (a)+(b) are PERSUADE-only; (c) changes the grade_level clauses in both corpora; (e) the
%             education arm is now a sex x ethnicity x economic-status factorial, so the sex and
%             ethnicity clauses are composite (all three attributes stated) and the ethnicity pole was
%             corrected from white to Black in the decision frame; and
%             (d) the rendered header now carries the assignment and the body keeps its paragraph
%             breaks, which changes every education item. The ASAP education numbers are not marked
%             \P individually: every one of them already carries \ddag, so none is quotable either way.
% A macro may carry SEVERAL marks (two independent reasons, each cleared by its own re-run), so the
% saved original is keyed by a counter rather than by the macro name: keying by name made the second
% \markstale of the same macro overwrite the first one's backup with the already-wrapped definition,
% i.e. define it in terms of itself (a TeX capacity error at first use, not a warning).
\newcounter{stalemark}
\newcommand{\markstale}[2]{%
  \stepcounter{stalemark}%
  \expandafter\let\csname stale@\thestalemark\expandafter\endcsname\csname #1\endcsname
  \expandafter\edef\csname #1\endcsname{%
    \noexpand\textcolor{orange}{%
      \expandafter\noexpand\csname stale@\thestalemark\endcsname
      \noexpand\textsuperscript{\unexpanded{#2}}}}}
% credit (invalid)
\markstale{autoInflIntersectionCreditbase}{\dag}
\markstale{autoInflIntersectionCreditnull}{\dag}
\markstale{autoInflSexCreditbase}{\dag}
\markstale{autoInflSexCreditnull}{\dag}
\markstale{crossAccBaselineCredit}{\dag}
\markstale{crossNCredit}{\dag}
\markstale{meanGapIntersectionCredit}{\dag}
\markstale{meanGapSexCredit}{\dag}
\markstale{additivityCosCredit}{\dag}
% ASAP (minor)
\markstale{crossAccBaselineEduAsap}{\ddag}
\markstale{crossNEduAsap}{\ddag}
\markstale{eduEthnicityExplicitAsap}{\ddag}
\markstale{eduEthnicityProxyAsap}{\ddag}
\markstale{eduGradeExplicitAsap}{\ddag}
\markstale{eduGradeProxyAsap}{\ddag}
\markstale{eduGradeProxyAsapnull}{\ddag}
\markstale{eduSexExplicitAsap}{\ddag}
\markstale{eduSexExplicitAsapnull}{\ddag}
\markstale{eduSexProxyAsap}{\ddag}
\markstale{nEvalStandpointAsap}{\ddag}
\markstale{standpointGapClassAsap}{\ddag}
\markstale{standpointGapControlAsap}{\ddag}
\markstale{standpointGapHobbyAsap}{\ddag}
\markstale{standpointGapIntersectionAsap}{\ddag}
\markstale{standpointGapOriginAsap}{\ddag}
\markstale{standpointGapPetAsap}{\ddag}
\markstale{standpointGapRaceAsap}{\ddag}
\markstale{standpointGapRegionAsap}{\ddag}
\markstale{standpointGapSexAsap}{\ddag}
\markstale{standpointMainEffectAsap}{\ddag}
% name lists / verdict wording / hiring factorial design
\markstale{batteryFamilyExplicitCV}{\S}
\markstale{batteryFamilyProxyCV}{\S}
\markstale{autoInflIntersectionCVbase}{\S}
\markstale{autoInflIntersectionCVnull}{\S}
\markstale{meanGapIntersectionCV}{\S}
\markstale{additivityCosCV}{\S}
\markstale{batterySexProxyCV}{\S}
\markstale{eduSexProxyAsap}{\S}
\markstale{eduSexProxyPersuade}{\S}
\markstale{eduEthnicityProxyAsap}{\S}
\markstale{eduEthnicityProxyPersuade}{\S}
\markstale{decisionDiscWinSex}{\S}
\markstale{decisionDiscWinAge}{\S}
\markstale{decisionDiscWinFamily}{\S}
\markstale{decisionDiscWinIntersection}{\S}
\markstale{decisionGapFairDiscSex}{\S}
\markstale{decisionGapFairDiscAge}{\S}
\markstale{decisionGapFairDiscFamily}{\S}
\markstale{decisionGapFairDiscIntersection}{\S}
% education substrate: stage poles widened (both corpora)
\markstale{eduGradeExplicitAsap}{\P}
\markstale{eduGradeProxyAsap}{\P}
\markstale{eduGradeProxyAsapnull}{\P}
% education substrate: PERSUADE record ids / sample / prompt selection changed
\markstale{crossAccBaselineEduPersuade}{\P}
\markstale{crossNEduPersuade}{\P}
\markstale{eduEthnicityExplicitPersuade}{\P}
\markstale{eduEthnicityProxyPersuade}{\P}
\markstale{eduGradeExplicitPersuade}{\P}
\markstale{eduGradeProxyPersuade}{\P}
\markstale{eduSexExplicitPersuade}{\P}
\markstale{eduSexProxyPersuade}{\P}
\markstale{nEvalStandpointPersuade}{\P}
\markstale{standpointGapClassPersuade}{\P}
\markstale{standpointGapControlPersuade}{\P}
\markstale{standpointGapHobbyPersuade}{\P}
\markstale{standpointGapIntersectionPersuade}{\P}
\markstale{standpointGapOriginPersuade}{\P}
\markstale{standpointGapPetPersuade}{\P}
\markstale{standpointGapRacePersuade}{\P}
\markstale{standpointGapRegionPersuade}{\P}
\markstale{standpointGapSexPersuade}{\P}
\markstale{standpointMainEffectPersuade}{\P}
  % hand-maintained / not-yet-serialized

% Keeps packages, notation, and the single-source-of-truth numbers consistent everywhere.
%
% NOTE: the final paper will later swap \documentclass{article} for the venue template (acmart for
% FAccT). This preamble stays largely reusable; drop the package lines the template already provides.

\usepackage[utf8]{inputenc}
\usepackage[T1]{fontenc}
\usepackage[margin=1in]{geometry}
\usepackage{amsmath,amssymb}
\usepackage{graphicx}
\usepackage{booktabs}
\usepackage[numbers]{natbib}  % numeric mode (plainnat); final paper swaps to the venue template's citations
\usepackage{xcolor}
\usepackage{hyperref}

% --- notation / shorthands ---------------------------------------------------------------------
\newcommand{\rewardmodel}{reward model}
\newcommand{\RMs}{reward models}
\newcommand{\diffmean}{difference-of-means}
\newcommand{\nullspace}{null-space projection}
\newcommand{\autoinf}{auto-influence}
\newcommand{\crossinf}{cross-influence}
\newcommand{\leace}{\textsc{Leace}}
\newcommand{\taco}{\textsc{TaCo}}
\newcommand{\todo}[1]{\textcolor{red}{[\textbf{TODO}: #1]}}

% --- single source of truth for numbers (see export_paper_numbers.py in the code repo) ----------
% AUTO-GENERATED by experiments/export_paper_numbers.py -- DO NOT EDIT BY HAND.
% source: /Users/CedricESMT/AI_safety/AI_Safety_Frankfurt/greenTeam/OneBiasAfterAnotherFork/artifacts/results/demographic/*.json | generated: 2026-07-02 | code commit: 63be6dc
% Hand-maintained / not-yet-serialized numbers live in shared/numbers_manual.tex.

\newcommand{\autoInflIntersectionCVbase}{1.00}
\newcommand{\autoInflIntersectionCVnull}{0.01}
\newcommand{\autoInflIntersectionCreditbase}{1.00}
\newcommand{\autoInflIntersectionCreditnull}{0.00}
\newcommand{\autoInflSexCVbase}{1.00}
\newcommand{\autoInflSexCVnull}{0.11}
\newcommand{\autoInflSexCreditbase}{1.00}
\newcommand{\autoInflSexCreditnull}{0.06}
\newcommand{\batteryAgeExplicitCV}{0.23}
\newcommand{\batteryAgeProxyCV}{0.99}
\newcommand{\batteryFamilyExplicitCV}{0.04}
\newcommand{\batteryFamilyProxyCV}{0.91}
\newcommand{\batterySexProxyCV}{0.63}
\newcommand{\cosCorrectnessPLvsCommute}{0.94}
\newcommand{\crossAccBaselineCV}{0.59}
\newcommand{\crossAccBaselineCredit}{0.46}
\newcommand{\decisionDiscWinAge}{0.00}
\newcommand{\decisionDiscWinFamily}{0.39}
\newcommand{\decisionDiscWinIntersection}{0.32}
\newcommand{\decisionDiscWinSex}{0.00}
\newcommand{\decisionGapFairDiscAge}{+2.51}
\newcommand{\decisionGapFairDiscFamily}{+0.27}
\newcommand{\decisionGapFairDiscIntersection}{+0.10}
\newcommand{\decisionGapFairDiscSex}{+4.92}
\newcommand{\erasureLinearLeaceCorrPL}{52\%}
\newcommand{\erasureLinearNoneCorrPL}{100\%}
\newcommand{\erasureMlpLeaceCorrPL}{100\%}
\newcommand{\meanGapIntersectionCV}{+0.62}
\newcommand{\meanGapIntersectionCredit}{+0.75}
\newcommand{\meanGapSexCV}{+0.59}
\newcommand{\meanGapSexCredit}{+0.39}
\newcommand{\modelSmall}{Skywork/Skywork-Reward-V2-Qwen3-0.6B}
\newcommand{\nEval}{200}
\newcommand{\probeConclAccHeldoutCommute}{100\%}
\newcommand{\probeConclAccHeldoutPL}{100\%}
\newcommand{\probeCorrAccHeldoutCommute}{100\%}
\newcommand{\probeCorrAccHeldoutPL}{100\%}
\newcommand{\reasonConclusionCommute}{+1.46}
\newcommand{\reasonConclusionIntersection}{+0.92}
\newcommand{\reasonConclusionPL}{+0.67}
\newcommand{\reasonCorrectnessCommute}{+1.46}
\newcommand{\reasonCorrectnessIntersection}{+1.33}
\newcommand{\reasonCorrectnessPL}{+1.45}
\newcommand{\transferCommuteToPLbase}{+1.45}
\newcommand{\transferCommuteToPLnull}{+0.17}
         % AUTO-GENERATED from results JSONs
% HAND-MAINTAINED numbers: values not yet serialized to JSON, or human-derived.
% Everything auto-generated lives in shared/numbers.tex (do not edit that one by hand).
%
% Additivity cosines: additivity_check.py now supports --out; until it is re-run with --out these
% stay here. Replace with the auto macros (\additivityCosCV / \additivityCosCredit) once serialized.
\newcommand{\additivityCosCV}{0.66}     % TODO auto: run_additivity.py --domain cv --out ...
\newcommand{\additivityCosCredit}{0.71} % TODO auto: run_additivity.py --domain credit --out ... (STALE, see below)

% --- STALE AUTO-NUMBERS (2026-09-16) ---------------------------------------------------------------
% numbers.tex is auto-generated and must not be edited, so stale values are flagged here instead: each
% \markstale{name}{mark} re-typesets the macro in orange with a mark. Every use in every document is
% flagged at once. DELETE THE MATCHING LINE once export_paper_numbers.py has re-exported the value from
% a re-run (a leftover line would flag a fresh number as stale).
%   \dag      credit: German Credit was decoded with the wrong UCI code table, and the pairs now come
%             from the sex x age x marital-status factorial -- every credit result must be re-run.
%   \ddag     ASAP: redaction-tag cleaning changed 22 of 12,976 essays (adjacent tags such as
%             @PERCENT1@NUM1) -- expected to be negligible, re-run for consistency.
%   \S        inputs changed: the proxy first names are now Haim et al.'s Table 3 (every name-proxy
%             number), and the decision-response verdicts now open without the synthetic "N years" field.
%   \P        education substrate changed: (a) the essay-id lookup keyed on the lossy `essay_id` column
%             and merged two essays, it now keys on `essay_id_comp`, so every PERSUADE record id and the
%             sampled set change; (b) the essays are now restricted to the stage-neutral prompts without
%             in-essay pupil cues; (c) the grade_level poles widened to 6th grade vs doctoral candidate.
%             (a)+(b) are PERSUADE-only; (c) changes the grade_level clauses in both corpora; (e) the
%             education arm is now a sex x ethnicity x economic-status factorial, so the sex and
%             ethnicity clauses are composite (all three attributes stated) and the ethnicity pole was
%             corrected from white to Black in the decision frame; and
%             (d) the rendered header now carries the assignment and the body keeps its paragraph
%             breaks, which changes every education item. The ASAP education numbers are not marked
%             \P individually: every one of them already carries \ddag, so none is quotable either way.
% A macro may carry SEVERAL marks (two independent reasons, each cleared by its own re-run), so the
% saved original is keyed by a counter rather than by the macro name: keying by name made the second
% \markstale of the same macro overwrite the first one's backup with the already-wrapped definition,
% i.e. define it in terms of itself (a TeX capacity error at first use, not a warning).
\newcounter{stalemark}
\newcommand{\markstale}[2]{%
  \stepcounter{stalemark}%
  \expandafter\let\csname stale@\thestalemark\expandafter\endcsname\csname #1\endcsname
  \expandafter\edef\csname #1\endcsname{%
    \noexpand\textcolor{orange}{%
      \expandafter\noexpand\csname stale@\thestalemark\endcsname
      \noexpand\textsuperscript{\unexpanded{#2}}}}}
% credit (invalid)
\markstale{autoInflIntersectionCreditbase}{\dag}
\markstale{autoInflIntersectionCreditnull}{\dag}
\markstale{autoInflSexCreditbase}{\dag}
\markstale{autoInflSexCreditnull}{\dag}
\markstale{crossAccBaselineCredit}{\dag}
\markstale{crossNCredit}{\dag}
\markstale{meanGapIntersectionCredit}{\dag}
\markstale{meanGapSexCredit}{\dag}
\markstale{additivityCosCredit}{\dag}
% ASAP (minor)
\markstale{crossAccBaselineEduAsap}{\ddag}
\markstale{crossNEduAsap}{\ddag}
\markstale{eduEthnicityExplicitAsap}{\ddag}
\markstale{eduEthnicityProxyAsap}{\ddag}
\markstale{eduGradeExplicitAsap}{\ddag}
\markstale{eduGradeProxyAsap}{\ddag}
\markstale{eduGradeProxyAsapnull}{\ddag}
\markstale{eduSexExplicitAsap}{\ddag}
\markstale{eduSexExplicitAsapnull}{\ddag}
\markstale{eduSexProxyAsap}{\ddag}
\markstale{nEvalStandpointAsap}{\ddag}
\markstale{standpointGapClassAsap}{\ddag}
\markstale{standpointGapControlAsap}{\ddag}
\markstale{standpointGapHobbyAsap}{\ddag}
\markstale{standpointGapIntersectionAsap}{\ddag}
\markstale{standpointGapOriginAsap}{\ddag}
\markstale{standpointGapPetAsap}{\ddag}
\markstale{standpointGapRaceAsap}{\ddag}
\markstale{standpointGapRegionAsap}{\ddag}
\markstale{standpointGapSexAsap}{\ddag}
\markstale{standpointMainEffectAsap}{\ddag}
% name lists / verdict wording / hiring factorial design
\markstale{batteryFamilyExplicitCV}{\S}
\markstale{batteryFamilyProxyCV}{\S}
\markstale{autoInflIntersectionCVbase}{\S}
\markstale{autoInflIntersectionCVnull}{\S}
\markstale{meanGapIntersectionCV}{\S}
\markstale{additivityCosCV}{\S}
\markstale{batterySexProxyCV}{\S}
\markstale{eduSexProxyAsap}{\S}
\markstale{eduSexProxyPersuade}{\S}
\markstale{eduEthnicityProxyAsap}{\S}
\markstale{eduEthnicityProxyPersuade}{\S}
\markstale{decisionDiscWinSex}{\S}
\markstale{decisionDiscWinAge}{\S}
\markstale{decisionDiscWinFamily}{\S}
\markstale{decisionDiscWinIntersection}{\S}
\markstale{decisionGapFairDiscSex}{\S}
\markstale{decisionGapFairDiscAge}{\S}
\markstale{decisionGapFairDiscFamily}{\S}
\markstale{decisionGapFairDiscIntersection}{\S}
% education substrate: stage poles widened (both corpora)
\markstale{eduGradeExplicitAsap}{\P}
\markstale{eduGradeProxyAsap}{\P}
\markstale{eduGradeProxyAsapnull}{\P}
% education substrate: PERSUADE record ids / sample / prompt selection changed
\markstale{crossAccBaselineEduPersuade}{\P}
\markstale{crossNEduPersuade}{\P}
\markstale{eduEthnicityExplicitPersuade}{\P}
\markstale{eduEthnicityProxyPersuade}{\P}
\markstale{eduGradeExplicitPersuade}{\P}
\markstale{eduGradeProxyPersuade}{\P}
\markstale{eduSexExplicitPersuade}{\P}
\markstale{eduSexProxyPersuade}{\P}
\markstale{nEvalStandpointPersuade}{\P}
\markstale{standpointGapClassPersuade}{\P}
\markstale{standpointGapControlPersuade}{\P}
\markstale{standpointGapHobbyPersuade}{\P}
\markstale{standpointGapIntersectionPersuade}{\P}
\markstale{standpointGapOriginPersuade}{\P}
\markstale{standpointGapPetPersuade}{\P}
\markstale{standpointGapRacePersuade}{\P}
\markstale{standpointGapRegionPersuade}{\P}
\markstale{standpointGapSexPersuade}{\P}
\markstale{standpointMainEffectPersuade}{\P}
  % hand-maintained / not-yet-serialized

\title{Working Notes --- One Judge After Another}
\author{Frankfurt AI Safety --- RM-Bias Working Group}
\date{Living document --- last built \today}
\begin{document}
\maketitle
\tableofcontents

\section{Research questions (draft)}
\label{sec:rq-draft}
Draft RQ set for the extension. Not yet locked; promote to \texttt{sections/} once settled.

\paragraph{Main question.}
Do language \RMs{} inherit EU-protected demographic (and intersectional) biases; are these biases
low- or high-complexity in the model's activations; can mechanistic reward shaping remove them; and is
that removal \emph{robust} --- or does the bias only get cleaned from the reward scalar while remaining
entangled in the representation and leaking back under new conditions?

\subsection{RQ0 --- Replication, and re-testing the base paper's own complexity claims}
Two limbs. \textbf{(a)} Do the base-paper low-complexity interventions (length, uncertainty, position)
reproduce on the eleven \RMs{}, with RewardBench-2 non-inferiority \citep{fein2026onebias}? \textbf{(b)}
Apply \leace{} $+$ non-linear-probe (\taco{}) to the base paper's \emph{own} identified biases (the
``low-complexity'' length/uncertainty/position and the ``high-complexity'' sycophancy/model-style): does
\nullspace{} remove the concept from the \emph{representation}, or only its contribution to the reward
\emph{scalar}? I.e.\ re-run the base paper's low/high-complexity classification under the stronger
erasure/recoverability test, and check the reward-vs-representation gap for the spurious biases too ---
not just the demographic ones. (This makes our robustness check a contribution back to the base method.)

\subsection{RQ1 --- Prevalence: which protected attributes are inherited?}
Which EU-protected attributes does an \rewardmodel{} exhibit a measurable bias for when scoring
otherwise-identical high-risk decision content across three Annex-III domains --- recruitment
(CV screening), credit scoring, and education (grading student assignments/tests)? Measured on
matched-pair contrastive items under explicit and proxy encodings, via \autoinf{} (preference shift)
and \crossinf{}
(reliability harm) \citep{kumar2025prefix}, across the eleven \RMs{}. We report magnitude and removability,
not a betted direction.

\paragraph{RQ1.1 --- Intersectional.}
For the anchor intersection (sex $\times$ age $\times$ family-status; women $\approx 30$)
\outdated{credit now uses sex $\times$ age $\times$ \emph{marital status}, see log 2026-09-16}, is there a
\emph{distinct} intersectional bias, and is it \emph{non-additive} --- larger than, and not equal to, the
vector sum of the single-attribute directions?

\subsection{RQ2 --- Downstream harm (decision-response)}
When the \rewardmodel{} scores a \emph{decision response} rather than the applicant text --- a fair
recommendation vs a clearly discriminatory one for the same applicant --- which does it reward? I.e.\
would the RM reward a policy model that discriminates (the RLHF / best-of-N propagation story)? Headline:
win-rate of the discriminatory verdict over the fair one, with an explicit \textbf{evasion control}
(a non-committal answer must not out-score a substantive fair one). Already implemented. Since
2026-09-16 in all three domains, each with a binary decision: advance to interview (hiring), approve the
loan (credit), pass or fail (education).

\subsection{RQ3 --- Complexity: low vs high, and where it lives}
For each attribute and the intersection, is the bias low- or high-complexity, distinguishing two senses:
\begin{itemize}
  \item \textbf{reward-scalar level:} is its contribution to the scalar reward a single near-linear
        direction that \nullspace{} removes at near-zero cost to legitimate accuracy?
  \item \textbf{representation level:} after provably-optimal linear erasure (\leace{}
        \citep{belrose2023leace}), can a non-linear (MLP) probe still recover the concept? Linearly
        erasable but non-linearly recoverable $=$ the \taco{} signature of an entangled (high-complexity)
        bias \citep{taco2024}. This upgrades the base paper's nulling-based classification.
\end{itemize}

\subsection{RQ4 --- Mitigation via mechanistic reward shaping}
Where the bias is (reward-scalar) low-complexity, does \diffmean{}/\leace{} + \nullspace{} reduce the
reliability harm and preference bias \emph{without} degrading legitimate reward accuracy (RewardBench-2
non-inferiority)? Report the trade-off curve over projection strength $\alpha$.

\subsection{RQ5 --- Robustness / generalization of the intervention}
Because reward shaping removes the bias's linear contribution to the scalar but not the entangled concept
from the representation (RQ3), how robust is the mitigation? Does a projection fit on one setting stay
effective on held-out domains, templates, encodings (explicit$\leftrightarrow$proxy), phrasings, and
unseen intersections --- or does the bias concept \emph{leak back} into the reward scalar under
distribution shift? (Operationalized: fit on split/domain A, measure residual bias $+$ accuracy on
held-out split/domain B; cross-domain and cross-encoding transfer.) The three domains --- credit,
recruitment (CV), and education (grading) --- give us a genuine cross-domain transfer test.

\subsection{RQ6 --- Cross-model consistency: dataset- vs architecture-driven}
Is the bias --- magnitude and low/high-complexity class --- consistent across the eleven \RMs{}, or does it
vary by model? Our set is a natural experiment with the controlled comparisons built in: the Skywork-V2
models share Skywork's preference data but differ in base architecture/size (Llama-3.1-8B, Qwen3-8B,
Qwen3-0.6B, Gemma-2-27B), the AllenAI RB2 pair holds the recipe fixed across 8B and 70B (a clean size
comparison), the two Nemotron-32B models hold the recipe fixed across base generations, and
OpenAssistant-DeBERTa (encoder) plus QRM-Gemma (quantile/distributional) vary the RM \emph{type} itself.
If the bias clusters by \emph{training data} it is
dataset-driven; if by \emph{architecture/family}, architecture-driven \citep{kumar2025prefix}. This
locates where mitigation must live (upstream data governance vs model-internal editing).

\subsection{Validity confounds to pre-empt}
\begin{itemize}
  \item \textbf{Reasoning-correctness (construct validity).} Is the measured effect really the protected
        attribute, or the RM tracking (correct) reasoning / response quality? We have evidence it is the
        latter for verdicts (reasoning-flip 2$\times$2 with a non-demographic commute control). Needs to
        be a stated control, not a footnote.
  \item \textbf{Brittleness.} RMs react to non-demographic perturbations too; matched pairs must hold all
        non-target variation fixed (our Tier-1 gate).
  \item \textbf{Evasion.} Fairness framing can manifest as non-answers; score for evasion so it does not
        masquerade as fairness.
\end{itemize}

\subsection{Deferred / potential future directions}
\begin{itemize}
  \item Group-level (demographic-parity) framing \citep{song2025groupfairness} as a complement to the
        same-content counterfactual design --- do the two fall on opposite sides of the complexity
        boundary? \emph{Deferred}: out of scope for now, since it needs a \emph{new data pipeline} (group
        construction from genuinely different prompts, e.g.\ arXiv-style occupational proxies) rather than
        our matched-pair generator.
\end{itemize}

\subsection{Design axes (factors, not RQs)}
These modulate RQ1--RQ4 rather than being separate questions:
domain $\{$credit, recruitment (CV), education (grading)$\}$; task framing
$\{$economic, fairness, neutral$\}$; encoding $\{$explicit, proxy$\}$; scoring readout
$\{$scalar matched-pair (primary), forced-choice pairwise (secondary cross-check)$\}$.

\subsection{Protected-attribute scope}
Ground the attribute list in EU Charter Art.~21 $+$ AI Act Art.~10(2)(f,g), Art.~15 \citep{euaiact2024}.
Anchors are per-domain: credit and recruitment (CV) anchor on the sex $\times$ age $\times$ family-status
intersection (the pregnancy window) \outdated{credit: sex $\times$ age $\times$ marital status as a full
factorial, since family-status markers contradict the credit profile's employment field}; \textbf{education (grading)} anchors instead on a header
\textbf{name} ($\to$ sex $+$ ethnicity name-proxy) and \textbf{grade level} ($\to$ age / education-stage),
injected on real essays; disability is deferred. All are drawn from the same scope below.
\begin{itemize}
  \item \textbf{Explicitly statable --- cleanest matched pairs:} sex, age, family/parental status,
        \textbf{disability}. Each can be a single-slot stated marker (``a woman'', ``50 years old'',
        ``on parental leave'', ``has a disability'') plus a proxy form --- a true single-axis swap.
  \item \textbf{Proxy-only (name-based) --- standard but confounded:} ethnic/national origin. There is
        \emph{no} clean stated form on a CV/loan profile; the audit-study method swaps validated
        \emph{name} lists coded to different origins \citep{haim2024whatsinaname}. Caveat: names conflate
        ethnicity with class, region, and religion, so the single-axis guarantee is weaker than for
        sex/age --- report it as a name-proxy effect, not pure ethnicity.
  \item \textbf{Hard to operationalize without artifice --- exploratory:} religion, sexual orientation.
        Explicit markers are unnatural on these profiles and proxies are noisy.
\end{itemize}

\section{Reward models under evaluation}
The reward-model (RM) set is chosen to disentangle whether bias is \emph{dataset-driven} or
\emph{architecture-driven}, with controlled comparisons built in: the same Skywork-V2 training data across
different base architectures (Llama-8B vs Qwen3-8B); the same training recipe across a large size gap (the
AllenAI RB2 8B vs 70B); distinct RM \emph{types} (a DeBERTa encoder RM and a quantile/distributional QRM);
and a frontier 70B flagship. \textbf{All reported runs go on the hessian.AI cluster (A100/H100 80\,GB),}
under a two-month allocation secured 2026-08. The Apple Silicon (MLX) path is retained for local prototyping
and as the numerical parity reference, not as a source of published numbers. \outdated{The MLX backend has
since been removed from the code repository; CUDA is the only path.} Everything routes through the
same probe / null-space / LEACE pipeline, unchanged per model. The \textbf{GPUs} column gives the bf16
scheduling requirement: everything up to 32B fits on a single 80\,GB card; the two 70B models need
$\approx$140\,GB and are sharded across two.
\begin{center}\footnotesize\setlength{\tabcolsep}{4pt}
\begin{tabular}{@{}llllp{4.1cm}@{}}
\hline
\textbf{Reward model} & \textbf{Params} & \textbf{Base} & \textbf{GPUs} & \textbf{Role in the comparison}\\
\hline
Skywork-Reward-V2-Qwen3-0.6B        & 0.6B & Qwen3   & 1 & pilot; whole pipeline validated here\\
OpenAssistant deberta-v3-large-v2   & 0.4B & DeBERTa & 1 & encoder RM, pair-format; architecture outlier\\
Skywork-Reward-V2-Llama-3.1-8B      & 8B   & Llama-3.1 & 1 & top open RM on RewardBench-2\\
Skywork-Reward-V2-Qwen3-8B          & 8B   & Qwen3   & 1 & same Skywork data, different base (arch effect)\\
AllenAI Llama-3.1-8B-Instruct-RM-RB2& 8B   & Llama-3.1 & 1 & different recipe; 8B half of the scaling pair\\
Skywork-Reward-Gemma-2-27B          & 27B  & Gemma-2 & 1 & Skywork data on Gemma\\
QRM-Gemma-2-27B (nicolinho)         & 27B  & Gemma-2 & 1 & quantile/distributional RM (RM-type effect)\\
Qwen-2.5-Nemotron-32B-Reward        & 32B  & Qwen2.5 & 1 & Nemotron recipe\\
Qwen-3-Nemotron-32B-Reward          & 32B  & Qwen3   & 1 & Nemotron recipe, newer base\\
AllenAI Llama-3.1-70B-Instruct-RM-RB2& 70B & Llama-3.1 & 2 & same recipe as the AllenAI 8B (controlled scaling)\\
Llama-3.3-Nemotron-70B-Reward       & 70B  & Llama-3.3 & 2 & frontier flagship (state-of-the-art open 70B)\\
\hline
\end{tabular}
\end{center}
\emph{Eleven RMs spanning 0.6B--70B and the Llama / Qwen / Gemma / DeBERTa families. GPUs = 80\,GB A100/H100
cards needed per run at bf16. Exact Hugging Face paths for the NVIDIA Nemotron models are to be confirmed
against the current releases.}

\section{Open questions / things we are chasing}
\begin{itemize}
  \item \todo{Does the reasoning-correctness effect survive on a larger / instruct RM, and does
        \crossinf{} become interpretable (acc\_baseline $>0.6$)? Currently \crossAccBaselineCV{} on \modelSmall{}.}
        \emph{Partly resolved:} on the \textbf{education} domain \crossinf{} \emph{is} interpretable even on
        \modelSmall{} (acc\_baseline \crossAccBaselineEduAsap{} ASAP / \crossAccBaselineEduPersuade{} PERSUADE)
        --- essay quality is RM-trackable; credit/CV still need a bigger RM.
        \outdated{Credit's chance-level acc\_baseline was measured on profiles decoded with the wrong code
        table (log 2026-09-16); re-measure before concluding the RM cannot track loan risk.}
        \emph{Interpretation caveat (2026-09-17):} credit's ``strong'' label (\texttt{credit\_good}) is a
        real repayment \emph{outcome}, not a rating of the profile. Every applicant had already passed the
        bank's own checks, and bad credits are heavily oversampled (30\% in the data vs.\ about 5\% real
        prevalence, \citealp{groemping2019southgerman}). A near-chance acc\_baseline may therefore say as
        much about how hard the outcome is to predict from these fields as about the RM. Say so when
        reporting credit \crossinf{}.
  \item \todo{RQ5 robustness: does a projection fit on one setting leak back under distribution shift?}
        No transfer experiment has been run yet; the three domains give us the substrate for it.
  \item \textbf{\leace{} on the demographic directions.} The erasure/non-linear-probe test has so far only
        been run on the \emph{reasoning} concepts of the model-response arm (hiring domain). Running it on the
        demographic directions is what would establish the reward-vs-representation gap for protected
        attributes directly. Until then the abstract/discussion claim is scoped to the reasoning concept.
  \item \textbf{Accuracy guardrail not yet reported.} RQ0/RQ4 promise RewardBench-2 non-inferiority as the
        ``without degrading legitimate accuracy'' check; it has not been run. Note that on credit and CV the
        RM is at chance on quality (\crossAccBaselineCredit{} / \crossAccBaselineCV{}), so in-domain there is
        no accuracy signal to preserve --- at present the trade-off is only measurable on education.
        \outdated{credit: re-measure on the corrected data.}
  \item \textbf{Forced-choice pairwise readout.} Listed as the secondary readout in the design grammar but
        never implemented: it is out-of-distribution for a scalar RM and needs a generative judge, so it is
        deferred rather than dropped.
  \item \textbf{RQ0b --- \leace{} retrofit to the base paper's own biases} (length, position, uncertainty).
        Framed as a contribution back to the base method; not yet run.
\end{itemize}

\subsection*{Validity workstream (item 2 delivered 2026-08-12)}
Three checks on whether our \emph{instruments} measure what we claim. These bear on how much of the A1 marker
evidence survives, so they are scheduled ahead of further breadth.
\begin{enumerate}
  \item \textbf{Real vs.\ synthetic marker validity --- promote from footnote to finding.} On the one real-field
        run (German Credit \texttt{personal\_status\_sex}, $n{=}142$/group, probe acc $0.92$) the real-field
        direction is \emph{anti}-correlated with our injected-marker directions: $\cos = -0.216$ against the
        synthetic family-status direction and $-0.297$ against the synthetic sex direction, with a real-field
        reward gap of $\approx 0$. Plan: extend to PERSUADE~2.0's unused real race/sex/ELL labels and to
        further credit fields; report $\cos(\text{real}, \text{synthetic})$ per axis; treat a persistent
        negative or near-zero cosine as a result \emph{about the injected-marker method itself}.
        \outdated{Void (log 2026-09-16): the run decoded \texttt{personal\_status\_sex} with the wrong
        code table, so its ``single'' men were married/widowed men and its comparison group mixed divorced
        men with single women. Sex is not even recoverable for the largest code. The corrected run contrasts
        divorced/separated with married/widowed men ($n{=}47$ per group) against the synthetic
        \emph{marital-status} direction; not yet run.}
        \emph{Now cheaper:} item 2 brings a second corpus with a real sex label, so this is measurable on
        hiring as well as credit without any new annotation.
  \item \textbf{Bias-in-Bios real CVs --- \textsc{done}.} The hiring domain now loads Bias-in-Bios
        \citep{dearteaga2019bias} instead of the synthetic CV generator (4000 matched pairs, 8 axis
        $\times$ encoding cells, 0\% gate discard). Quality is \emph{role match} (true occupation $==$ the
        target role named in the header), so hiring gains a quality axis the RM can plausibly judge --- the
        thing credit lacks. \textbf{Caveat that matters:} the substrate changed, the evidence has not. The
        erasure, reasoning and decision arms have \emph{not} been re-run, so every hiring number in the
        write-up still rests on synthetic CVs. \textbf{Open sub-task:} the corpus preserves names and
        pronouns, so bodies are scrubbed; the Tier-1 gate structurally cannot verify that (leaked text sits
        on both poles and cancels), so \texttt{validate\_bios\_scrub.py} probes the corpus's \emph{real}
        gender label out of scrubbed activations. Near chance $\Rightarrow$ clean. That check has not been
        run on \modelSmall{} yet, and no sex-axis number from this substrate should be trusted until it is.
        \outdated{The check has since been run (review site: +0.10 over chance, about half of it the
        occupational prior), but on the scrub before 2026-09-17; re-run on the revised scrub.}
  \item \textbf{Name-pool validation} against Bertrand--Mullainathan / Haim et al.\ before any publishable
        ethnicity number. \texttt{markers.py} already self-flags the pools as an unvalidated starter set, and
        ethnicity is exactly the axis that is template-unstable across corpora
        (\eduEthnicityExplicitAsap{} ASAP vs \eduEthnicityExplicitPersuade{} PERSUADE).
        \emph{Done 2026-09-16:} the pools are now Haim et al.'s Table~3 (40 names). Remaining gap: the
        assumption that a name's gender signal dominates its age and class signals --- see the log entry
        of that date for the planned check.
\end{enumerate}

\subsection*{Codebase: prune and migrate to a fresh repository (planned)}
\outdated{Under way ahead of the schedule below: the code now lives in the derivative repository
\texttt{One-Judge-After-Another} (not a fork; upstream cited in \texttt{NOTICE}), with the MLX backend and
the upstream arms removed. The code is being reviewed folder by folder before it is committed.}
The work lives in a fork of \texttt{drfein/OneBiasAfterAnother}, and that has become the wrong container for
it:
\begin{itemize}
  \item \textbf{We do not reproduce the upstream results.} The fork carries the base paper's length,
        position, sycophancy and uncertainty arms, which we neither run nor report --- the extension is the
        whole contribution. Keeping them invites the reader to assume a replication we are not performing
        (RQ0 aside, which is a scoped re-test of their \emph{claims}, not a re-run of their pipeline).
  \item \textbf{The MLX backend has served its purpose.} It existed so the pipeline could be built and
        validated on Apple Silicon while we had no GPUs. With the allocation confirmed and CUDA parity
        verified (log, 2026-09-09), every reported number now comes off the cluster, so a second backend is
        maintenance cost plus a second numerical path to reason about.
  \item \textbf{We are not admins on the fork.} Discussions, Pages and app installs are unsettable there;
        the experiment-review site already had to live in a separate repo we own for exactly this reason.
\end{itemize}
\noindent \textbf{Plan:} copy the demographic pipeline, the site builder and the \texttt{cluster/} setup into
a fresh repository we own, keep the fork read-only as provenance, and \emph{cite} the base paper rather than
vendoring it. Two things to carry over deliberately rather than let lapse: the \textbf{MLX parity evidence},
which is what licenses the claim that the CUDA numbers are sound even after the backend is dropped, and the
snapshot in \texttt{results/}, which is the only public record of the pilot numbers.
\todo{Schedule this \emph{after} the eleven-RM ladder completes. A repo move mid-run would invalidate every
path in \texttt{cluster/config.yaml}, the PFSS working directory and the Hugging Face cache layout.}

\section{Target venues}
Deadlines below are as of \textbf{2026-08-04}; re-check before relying on any of them. The plan is
\textbf{two-track}: put the pilot findings into a workshop \emph{now}, and hold the full eleven-RM scaling
paper for an archival venue once the compute lands. The same results support two framings, and the framing
should follow the venue --- an \emph{interpretability / method} framing (the \nullspace{} $+$ \leace{} $+$
non-linear-probe recoverability story, i.e.\ the reward-vs-representation gap) for ICLR/NeurIPS, and a
\emph{fairness / auditing} framing (protected and intersectional attributes, a released bias benchmark, a
practitioner-usable audit method) for FAccT/AIES.

\begin{center}\footnotesize\setlength{\tabcolsep}{4pt}
\begin{tabular}{@{}p{2.5cm}p{3.4cm}p{9.0cm}@{}}
\hline
\textbf{Venue} & \textbf{Next deadline} & \textbf{Fit / note}\\
\hline
\multicolumn{3}{@{}l}{\textbf{Tier 1 --- primary archival targets (the full eleven-RM scaling paper)}}\\
ICLR 2027 & abstract 2026-09-19, paper 2026-09-24 & Best near-term big-ML target: the mechanistic
contribution (projection $+$ erasure $+$ non-linear recoverability) is the kind ICLR rewards. About seven
weeks out --- realistic only if the scaling runs land.\\
ACL / EMNLP (via ARR) & rolling; the EMNLP 2026 May cycle has closed, so a later ARR cycle aimed at ACL 2027
& RM bias on real essays is squarely NLP, and ``Ethics, Bias, and Fairness'' is a named track. The rolling
cadence lets us submit when results are ready instead of racing one date.\\
ACM FAccT 2027 & $\approx$ Jan 2027 (2026 was abstract Jan 8 / paper Jan 13) & The premier fairness venue and
the closest match to our framing: protected $+$ intersectional attributes, a released benchmark, an auditing
method. Strong second archival submission with the fairness-forward spin.\\
\hline
\multicolumn{3}{@{}l}{\textbf{Tier 2 --- ethics/society specialist venues}}\\
AIES 2027 & $\approx$ May 2027 & Excellent topical fit. The 2026 cycle closed 2026-05-21 (conference
Malm\"o, Oct 2026), so this is a next-year target.\\
AAAI 2027 & closed 2026-07-27 & Deadline has passed; it would have been too early regardless, with the
scaling work still pending compute.\\
\hline
\multicolumn{3}{@{}l}{\textbf{Tier 3 --- workshops (the near-term move)}}\\
SoLaR & typically $\approx$ Sep; 2026 edition to confirm & Socially Responsible Language Modelling Research.
Near-perfect fit: its named topics are LM interpretability, alignment/safety, and bias/fairness. Ran at
NeurIPS 2023/24 and COLM 2025.\\
NeurIPS 2026 / ICLR 2027 interpretability workshops & with the host conference & E.g.\ ``Interpretability for
Discovery'' at NeurIPS 2026. Good homes for the mechanistic-erasure angle specifically.\\
\hline
\end{tabular}
\end{center}

\begin{itemize}
  \item \textbf{Why workshop first.} The write-up is explicitly in progress and the headline results still
        come from a single 0.6B RM. A workshop paper puts the reward-vs-representation gap and the
        standpoint-credibility effect in front of exactly the right community now, and de-risks the later
        archival submission.
  \item \textbf{September is the decision point.} ICLR 2027 (2026-09-24) and the likely SoLaR / NeurIPS
        workshop deadlines both fall around September 2026. Compute is no longer the constraint (the
        hessian.AI allocation runs 2026-08 to 2026-10); what decides archival-now vs.\ workshop-now is
        simply whether the eleven-RM runs are finished and analysed by then.
  \item \textbf{Sequencing.} Workshop paper on the pilot findings $\rightarrow$ ICLR 2027 or FAccT 2027 for
        the full study. FAccT's January deadline is the natural fallback if September is too tight.
  \item \todo{Confirm the SoLaR 2026 edition, its host conference and deadline; confirm FAccT 2027 dates
        when the CFP posts; confirm the ARR cycle that feeds ACL 2027.}
\end{itemize}

\section{Acknowledgement (required in the final submission)}
Compute for the scaling runs comes from the hessian.AI 42 cluster, and the allocation carries a
funding-acknowledgement obligation. The following must appear \textbf{verbatim} in any paper reporting these
results:

\begin{quote}
We gratefully acknowledge support from the hessian.AI Service Center (funded by the German Federal Ministry
of Research, Technology and Space, BMFTR, grant no.\ 16IS22091) and the hessian.AI Innovation Lab (funded by
the Hessian Ministry for Digital Strategy and Innovation, grant no.\ S-DIW04/0013/003).
\end{quote}

\noindent This is a condition of the allocation rather than a courtesy, so it belongs on the camera-ready
checklist for every venue in the section above.
\todo{Double-blind caveat: the text names funders and grant numbers and would deanonymise us. ICLR, FAccT
and AIES all review double-blind, so it must be \emph{withheld} from the submitted version and added at
camera-ready --- the opposite of the usual instinct to paste acknowledgements in early.}

\section{Log}
\subsection*{2026-07-02 --- domain scope \& direction decisions}
\begin{itemize}
  \item Team decided to add a \textbf{third domain: education (grading student assignments/tests)}, as a
        direct-scoring parity arm to credit + recruitment (CV). Graded object $=$ student answer quality
        (strong/weak $=$ the \crossinf{} quality label, grading analog of CV \texttt{qualified} / credit
        \texttt{credit\_good}); axes $=$ sex $+$ ethnicity (name-proxy) $+$ disability. \emph{Docs updated
        here; substrate/code not yet built --- results pending.}
  \item Team decided to make \textbf{intersectional protected-attribute bias in high-risk decisions the
        project's sole focus}; the other candidate tracks under consideration at the time (dialect/register,
        political content-vs-style, generational register) were closed out rather than run in parallel.
\end{itemize}
\subsection*{2026-07-04 --- education arm built + first results (\modelSmall{}, ASAP $+$ PERSUADE)}
Built the education arm on \emph{real} essay corpora (PERSUADE 2.0 $+$ ASAP-AES; approach A1 --- header
proxy, essay body held byte-identical) rather than synthetic answers. Axes: header name $\to$ sex /
ethnicity; grade level $\to$ age. First results on \modelSmall{}:
\begin{itemize}
  \item \textbf{Two large, low-complexity biases replicate across both corpora.} Explicit \textbf{sex}
        auto-influence 0.88 (ASAP) / 0.95 (PERSUADE) --- stating ``a woman/man'' --- but the name proxy is
        only 0.24 / 0.35 (for sex, explicit $\gg$ proxy). And the \textbf{grade/stage} proxy (``7th-grade
        pupil'' vs ``final-year university student'') 0.96 / 0.97 --- while \emph{stated} age is near-zero
        (0.37 / 0.03); the RM rewards the \emph{identical} essay far more when it is attributed to a
        university student. Both collapse under the $\alpha$-sweep ($\to \sim$0.1) $\Rightarrow$ low-complexity.
  \item \textbf{Ethnicity is weak and template-direction-unstable} on both corpora (the two templates prefer
        opposite poles, so the aggregate washes out) --- needs more templates/names before any claim.
  \item \textbf{\crossinf{} is now interpretable:} acc\_baseline 0.72 (ASAP) / 0.88 (PERSUADE) $>0.6$ ---
        essay quality is what the RM natively scores (vs credit 0.51 \outdated{wrong code table} / CV 0.59). But CI\_base $\approx$ CI\_ref
        on every axis $\Rightarrow$ a symmetric \emph{marker main-effect}, not a protected-attribute-specific
        reliability harm at 0.6B (a legitimate null to revisit at 8B).
\end{itemize}
Runner generalizations: \texttt{run\_experiment --dataset-source} now honored under \texttt{--config};
\texttt{run\_demographic\_battery}/\texttt{run\_crossinfluence} made domain-aware. Deferred: PERSUADE real
demographic labels / group-level; explicit-ethnicity template instability; disability axis; the 8B RM.

\subsection*{2026-07-04 --- positioned-argument (A2) arm: standpoint-credibility bias}
A2 makes identity \emph{load-bearing for the argument's authority}: a first-person positionality sentence
grounds a real PERSUADE essay's conclusion in the author's lived experience, and we swap only the claimed
identity (``as \{a woman / a man / a Black person / \dots / a working-class Black woman / an affluent white
man\} who has lived these realities firsthand, I am convinced this is the right conclusion''), holding the
essay body byte-identical. Injection position is a config parameter (v1 $=$ conclusion). We add a
\textbf{non-demographic control} family (gardener/cyclist, dog/cat, rural/city) and a \textbf{neutral
(no-positionality) baseline}, and decompose each axis into a \emph{main effect} (any standpoint vs neutral)
and a magnitude \emph{identity gap} ($\Delta_a-\Delta_b$). \textbf{Result (\modelSmall{}, PERSUADE,
$n=160$):}
\begin{itemize}
  \item \textbf{General sycophancy for lived experience:} the main effect is positive on every axis
        ($+0.14$ to $+0.38$) --- the RM rewards \emph{any} ``I have lived this'' appeal over the plain essay,
        even hobby/pet.
  \item \textbf{A genuine, demographic-specific pro-marginalized bias:} the identity gap is large on the
        demographic axes (sex $+0.20$, race $+0.22$, class $+0.36$, origin $+0.28$, \textbf{intersection
        $+0.70$}) while the \emph{clean neutral controls are $\approx 0$} (gardener-vs-cyclist $-0.01$,
        dog-vs-cat $+0.02$). Direction is pro-marginalized (woman$>$man, Black$>$white, working-class$>$
        wealthy, immigrant$>$native), strongest for the working-class Black woman.
  \item \textbf{Methodological caution:} auto-influence \emph{saturates} (direction-consistency, not
        magnitude --- a trivial $+0.09$ gap gave auto-influence $0.80$), and an authority-tinged control
        (``retired teacher'') is contaminated. \textbf{Report the identity-gap magnitude against clean
        controls, not auto-influence, for A2.}
\end{itemize}
Deferred: opening/middle/random positions; referee (policy-model-grades $\to$ RM-rewards) framing; ASAP
replication; social-issue corpus where standpoint is topic-intrinsic; 8B RM.

\subsection*{2026-07-05 --- A2 hardening: robustness + stance control}
The pro-marginalized identity gap survives four robustness checks (\modelSmall{}, $n\approx120$--$160$), with
the clean neutral controls staying $\approx 0$ throughout:
\begin{itemize}
  \item \textbf{Injection position:} holds at conclusion / opening / middle / random (intersection gap
        $0.71 / 1.20 / 0.48 / 0.49$). Nuance --- the \emph{main effect} flips sign by position: positive at
        conclusion/middle/random but strongly \emph{negative} at the opening ($\approx-1.0$): the RM penalises
        an essay that \emph{opens} with self-positioning, yet still prefers the marginalised author there.
  \item \textbf{Corpus:} replicates on ASAP (sex $0.25$, race $0.33$, class $0.47$, origin $0.49$,
        intersection $0.93$; controls $\le 0.09$).
  \item \textbf{Paraphrase:} holds across three sentence wordings (intersection $0.71 / 0.53 / 0.47$) --- not
        tied to one phrasing.
  \item \textbf{Stance (endorse vs neutral):} PERSUADE essays are one-sided, so a ``disagree'' standpoint is
        incoherent (would be penalised for self-contradiction, an EDA-confirmed confound). Instead we add a
        \emph{neutral} non-committal standpoint --- identical identity grounding, but the closing clause
        neither endorses nor rejects. The gap \emph{persists} under neutral (intersection $0.70\to0.74$, sex
        $0.20\to0.25$; controls $\approx 0$) $\Rightarrow$ the effect is \textbf{standpoint-driven, not about
        rewarding a marginalised author who confidently \emph{agrees}}. (Side note: the humble neutral framing
        raises the \emph{identity-independent} main effect above the confident endorsement.)
\end{itemize}
Tooling: \texttt{run\_positioned\_maineffect.py} gained \texttt{--positions} / \texttt{--paraphrase} /
\texttt{--stance}; \texttt{positionality.py} has paraphrase pools $+$ a neutral-stance variant. Deferred:
referee framing; a topic-natural social-issue corpus; 8B RM.
\subsection*{2026-08-12 --- compute secured; hiring substrate moved to real data}
\begin{itemize}
  \item \textbf{hessian.AI allocation:} two months of A100/H100 (80\,GB), 2026-08 to 2026-10. This covers the
        whole eleven-RM ladder, so the local-vs-rented split in the RM table above is retired --- all reported
        runs go on the cluster, and MLX stays as the prototyping / parity path. Supersedes the compute ask in
        \texttt{progress\_2026-07.tex} (kept as a July snapshot).
  \item \textbf{Hiring domain moved off the synthetic substrate} to \textbf{Bias-in-Bios} (De-Arteaga et al.\
        2019; \texttt{LabHC/bias\_in\_bios}, MIT): $\approx$400k real biographies, 28 occupations, with
        \emph{real} gender labels. Addresses validity item 2 --- our most load-bearing arms (erasure,
        reasoning, decision-response) had been sitting on entirely synthetic CVs.
  \item \textbf{Quality label = role-match.} The rendered header states a target role;
        \texttt{qualified} $=$ (true profession $==$ target role), assigned $\approx$50/50 at load time. This
        gives hiring a quality axis the RM can actually judge, which is exactly what credit and CV lack
        (\texttt{acc\_baseline} \crossAccBaselineCredit{} / \crossAccBaselineCV{}, at chance, which is why
        \crossinf{} is uninterpretable there).
  \item \textbf{Open risk:} the HF mirror ships only \texttt{hard\_text}, which \emph{preserves} names and
        pronouns, so we scrub ourselves. The Tier-1 gate cannot catch a bad scrub (it only compares the two
        poles to each other, so leaked gendered text is present on both sides and passes). Validated
        separately by probing the \emph{real} gender label from scrubbed bios --- near chance means the scrub
        held.
  \item The real gender labels also make validity item 1 measurable on a second domain: $\cos$(real-gender
        direction, synthetic sex-marker direction), which so far rests on one German Credit run that returned
        $-0.297$ \outdated{void, see 2026-09-16}.
\end{itemize}
\subsection*{2026-09-09 --- cluster access live; CUDA path validated}
\begin{itemize}
  \item \textbf{The hessian.AI 42 cluster is running our code.} Workspace \texttt{IL\_rm\_bias} on
        A100-SXM4-80GB via the MLDE/Determined platform. The setup is committed under \texttt{cluster/} in
        the code repo: task configs, environment notes, a checkpoint prefetcher and a runbook.
  \item \textbf{Parity smoke test passed.} The credit explicit-sex arm reproduces the Mac/MLX reference:
        baseline \autoinf{} exactly $1.00$, and --- the decisive line --- the nulled effect in raw reward
        units agrees to three decimals (\texttt{abs\_mean\_gap} $0.0689$ vs $0.069$) across two frameworks,
        two dtypes and two backends. \outdated{The parity conclusion stands (both backends scored the same
        data), but these values are no longer reference numbers: the credit data they were computed on has
        been replaced (2026-09-16). Re-establish the cluster reference after regenerating it.}
  \item \textbf{The single divergence is one we had already documented.} Nulled \autoinf{} reads $0.11$
        against a reference $0.06$. Nulling drives \texttt{mean\_gap} to $0.008$, i.e.\ essentially zero, so
        every A/B comparison is decided by numerical noise; \texttt{pref\_a\_rate} $0.555$ vs $0.53$ is
        $0.7$\,SE at $n{=}200$. This is the bf16-vs-fp32 decision-tie effect recorded during the original MLX
        parity work, not a framework fault.
  \item \textbf{Metric consequence, worth settling before eleven models of output exist.} \autoinf{} is a
        preference \emph{rate}: it saturates at $1.00$ for any consistently-signed effect and degenerates to
        noise once an effect is nulled --- and both ends of its range are exactly where we read it. Report
        \texttt{mean\_gap}/\texttt{abs\_mean\_gap} as the primary magnitude alongside it, as the A2 arm
        already does with \texttt{identity\_gap}.
  \item Consequence for the codebase: with CUDA validated, the MLX backend is no longer on the critical path
        --- see the migration plan above.
\end{itemize}

\subsection*{2026-09-16 --- substrates review: German Credit code table corrected, credit redesigned}
First line-by-line review session (\texttt{substrates/}). Every credit number in this document and the paper
is now stale (flagged \dag{} automatically via \texttt{shared/numbers\_manual.tex}); ASAP numbers carry a
minor flag (\ddag).
\begin{itemize}
  \item \textbf{The UCI code table for German Credit is wrong.} \citet{groemping2019southgerman} traced the
        data to its German sources and published a corrected table; the one shipped as \texttt{german.doc}
        (and copied by our decoder) mislabels checking account, credit history, purpose, savings, property,
        housing and personal status/sex, swaps foreign worker and people liable, and presents installment
        rate, number of credits and dependents as counts/percentages when they are bands. We verified the
        corrected table against our own \texttt{german.data}: the per-class frequencies match Grömping's
        Table~1 exactly (housing off by one record, a typo the report mentions); the old mapping does not.
        Consequences:
        \begin{itemize}
          \item every rendered credit profile described a different applicant from the one behind its
                good/bad label --- which may be part of why credit \crossinf{} sat at chance;
          \item \textbf{the real-field result is void} ($\cos=-0.297$ / $-0.216$, $n{=}142$): its ``single
                men'' were married/widowed men and its comparison group mixed divorced men with single
                women. Code A92 pools non-single women with single men, so sex is unknown for 31\% of
                applicants. The runner now contrasts divorced/separated with married/widowed men ($n{=}47$
                per group) against the synthetic marital-status direction;
          \item all credit pairs and results must be regenerated, and the cluster parity reference
                re-established.
        \end{itemize}
  \item \textbf{Consistency filter.} Records whose rendered fields contradict each other are dropped
        (unemployed with a skilled/management job: 45; owned housing without real estate: 4), leaving 952.
  \item \textbf{Credit is now a sex $\times$ age $\times$ marital-status full factorial.} Each applicant is
        rendered in all 8 cells (``The applicant is a 30-year-old married woman.''); single-attribute pairs
        flip one attribute with the other two stated and equal, balanced over their levels; the intersection
        pair is (female, 30, married) vs.\ (male, 50, single). Proxy: first name for sex, birth year for age;
        marital status has no clean proxy and stays explicit (no proxy marital pairs). 35{,}332 pairs, no
        gate discards.
        \begin{itemize}
          \item \textbf{Family status replaced by marital status for credit:} the parental-leave / career-gap
                markers contradicted the employment field (62 and 427 records), and only one side of each
                pair clashed, so the RM could punish implausibility rather than the attribute. Hiring keeps
                family status.
          \item \textbf{Plausibility filter:} ``single'' next to three or more dependents is unusual while
                ``married'' is not, so those 149 records are dropped (803 remain). \textbf{Reporting caveat:}
                130 of them are married/widowed men, whose share falls from 54.5\% to 48.4\%.
        \end{itemize}
  \item \textbf{Split by applicant.} The factorial yields several pairs per applicant, so probe/eval are now
        split by record (credit only; hiring and education still split per pair and can leak across
        templates --- flagged for the \texttt{scoring/} review). \texttt{probe\_size: 300} now means only
        $\approx$38 applicants per single-attribute direction; the value will be raised after timing runs on
        the cluster.
  \item \textbf{Reproducibility:} the credit generator seeded its RNG with Python's per-process-salted
        \texttt{hash()}, so earlier credit datasets were not exactly reproducible; fixed. The hiring and
        education generators still do this (\texttt{runners/} review).
  \item \textbf{Smaller fixes:} ASAP redaction tags (\texttt{@PERSON@}, adjacent \texttt{@PERCENT1@NUM1}) are
        now removed cleanly (22 of 12{,}976 essays change, hence the \ddag{} flag); a PERSUADE file without an
        essay-id column would have yielded $\sim$11 copies per essay (ours has one, no change); ASAP sets whose
        score cutoffs coincide are skipped instead of being labelled all-weak (no such set in our file); the
        job labels drop ``(non-)resident'', a nationality cue.
  \item \textbf{Name pools replaced with \citet{haim2024whatsinaname}, Table~3.} 10 first names per race
        $\times$ gender group, chosen from Gaddis's perception data for the most consistent racial
        perception (the paper adds the surnames Becker / Washington; we use first names only). The old
        starter pools overlapped only partly (white female 4 of 6, white male 1 of 6, Black female 4 of 6,
        Black male 3 of 6). Every name-proxy number is now stale (\S): CV proxy sex, education proxy sex
        and ethnicity, and the credit proxy cells.
        \begin{itemize}
          \item \textbf{Assumption (not validated):} a name's \emph{gender} signal is much larger than its
                \emph{age-cohort} and \emph{class} signals. Haim et al.\ validate perceived race and do not
                discuss age or class, and several names are cohort-marked (Katelyn, Hunter read younger than
                Laurie, Todd). If the assumption fails, a proxy-sex effect partly measures age or class.
                In the credit factorial the age pairs hold the name fixed, so there only the sex contrast is
                exposed.
          \item \textbf{Planned check (leave-names-out).} (a) Split each pool into folds, build the sex
                direction from pairs using the training names and evaluate on pairs whose names are held
                out; a stable effect across folds means no small set of names carries it. (b) Estimate the
                effect per name and regress it on a cohort score (e.g.\ the median birth year of each name
                in U.S. Social Security baby-name counts) and, where available, a class or education
                perception score; slopes near zero support the assumption. (c) Optionally, repeat with a
                second, age-matched name list. \todo{Schedule after the credit and education re-runs.}
        \end{itemize}
  \item \textbf{Decision-response arm extended to all three domains.} Each domain now asks its own binary
        question after the marked profile --- hiring: advance to an interview; credit: approve the loan;
        education: should the essay pass (pass/fail keeps it binary). Verdicts share an opening that
        concedes the merits and no longer read the synthetic ``N years of experience'' field (which had
        silently read ``several years'' on real biographies); the runner uses only strong records, so the
        positive decision is the correct one. Credit axes are sex, age, marital status and the intersection
        (no proxy for marital status); education axes are sex, ethnicity (Black-coded pole targeted) and
        grade level. New configs for credit and education. The reasoning 2$\times$2 is still hiring-only
        and still reads the years field. The existing hiring decision numbers are stale (\S).
\end{itemize}

\subsection*{2026-09-17 --- substrates review: credit files done}
\begin{itemize}
  \item Line-by-line review of the credit part of \texttt{substrates/} finished (\texttt{credit\_ingest},
        \texttt{credit\_clean}, \texttt{credit\_render}, the credit entry of \texttt{domains}). No
        correctness problems beyond those fixed on 2026-09-16.
  \item The rule filter keeps the good/bad balance: 564 of the 803 retained applicants are good credit
        (70.2\% vs.\ 70.0\% in the full data).
  \item \textbf{Label caveat:} \texttt{credit\_good} is an oversampled repayment outcome among applicants
        the bank had already accepted, not a rating of the profile (see the open question on
        acc\_baseline above).
  \item Small fixes: \texttt{run\_crossinfluence --raw-path} was silently ignored for credit and is now
        honoured (same filtered population); credit's axis list in the domain registry is now derived from
        the factorial instead of being a second copy.
  \item Relative data paths are fine as long as runs start from the repository root, which the cluster
        task config guarantees (\texttt{work\_dir} is the repository).
  \item \textbf{Hiring moved to a full factorial too} (sex $\times$ age $\times$ family status), for the
        same reason as credit: the pregnancy-window question is an interaction question, and the old
        design could only compare single-attribute pairs against one corner pair, with a direction
        cosine (\additivityCosCV{}) as the only interaction evidence. All 8 cells are now rendered per
        biography, so the motherhood interaction is a direct contrast.
        \begin{itemize}
          \item \textbf{The family proxy changed from a career gap to a parent-association role}
                (adapting \citet{correll2007motherhood}). A career break is not specific to parenthood
                and contradicts bios describing a continuing career, and only on one side of the pair.
                So the earlier CV finding ``family proxy \batteryFamilyProxyCV{} $\gg$ explicit
                \batteryFamilyExplicitCV{}'' may have been a career-gap penalty, not a family-status one.
                The proxy now measures \emph{parenthood} while the explicit clause states \emph{parental
                leave}; the labels say which.
          \item \textbf{Plausibility rules} drop bios whose content contradicts a stated age of 30
                (a year before 2014, $>$8 years of experience, ``two decades'', retirement):
                94{,}226 scrubbed bios leave 59{,}838, far more than the $\approx$4{,}000 sampled.
          \item Hiring pairs are now grouped by biography in the probe/eval split (as credit), which also
                closes the leak flagged for this domain; \texttt{generate\_bios.py} takes
                \texttt{--n-bios} instead of \texttt{--n-per} and writes \texttt{cells.jsonl}.
          \item Education keeps its single-axis design (its ``grade level'' axis is itself an age cue;
                a factorial there needs its own design discussion).
        \end{itemize}
  \item \textbf{Bias-in-Bios scrub revised} (checked on 20k and then all 257k bios).
        \begin{itemize}
          \item \textbf{The name check was mostly a no-op.} The corpus text has its first sentence removed,
                so most bios open with ``He''/``She'', and the scrub took that pronoun for the subject's
                name: for 82\% of kept bios no real name was ever identified, and 36\% of otherwise clean
                bodies still named a person, mostly the subject (``\dots John holds a Juris Doctorate'').
          \item \textbf{Fix: drop every bio that still names a person by a sex-coded first name.} List of
                1{,}767 U.S. first names from Social Security baby-name data (births 1940--2000, $\geq$90\%
                one sex; ordinary words, calendar and place names removed), built by
                \texttt{runners/build\_first\_names.py}; place and institution contexts are skipped.
          \item Also: titles are removed only before a name (the old pattern also deleted ``MS'' and
                ``miss''); only the full name and first name are replaced, surnames kept (they mangled text
                such as ``Young Adults''); emails and URLs replaced; more gendered person nouns trigger a
                drop (\emph{actress}, \emph{chairman}, \emph{brother}, \dots).
          \item \textbf{Result:} 94{,}226 of 257{,}478 bios kept (dropped: 98{,}332 length, 11{,}382 gendered
                words, 53{,}538 first names, mixed-case names included); the old ``no name found'' drop is gone. Women are 46.1\% of the
                corpus but 42.9\% of kept bios --- a small selection effect to report. ``women''/``men'' are
                kept and reported: 3.9\% of kept bios mention them.
          \item Every hiring number must be re-run on the new data anyway, and the scrub check
                (\texttt{validate\_bios\_scrub.py}) must be re-run before any sex-axis result is trusted.
        \end{itemize}
  \item \textbf{Role-match label leaked through the role name --- fixed.} Unqualified bios used to get a
        uniformly drawn other profession as target role, while qualified ones kept their own, so common
        professions appeared as targets mostly qualified and rare ones mostly not (``professor'' 95\%
        qualified, ``DJ'', ``rapper'', ``personal trainer'' 6\%). The role name alone predicted the label
        with \textbf{76.5\%} accuracy, so a hiring \texttt{acc\_baseline} could not tell ``reads the
        bio'' from ``likes professor profiles''. Now each bio's coin is fixed by (seed, row) and unqualified
        bios receive another unqualified bio's profession (shuffled, never their own): role-name-only
        accuracy \textbf{50.6\%}, P(qualified $\mid$ role) between 0.47 and 0.54 for every role. Any earlier
        hiring result that uses the \texttt{qualified} label predates this.
  \item Smaller \texttt{bios\_ingest} fixes: the unscrubbed text is kept in memory only when the scrub
        check asks for it; unknown professions raise; ``DJ'' is capitalised in the header.
  \item \textbf{Education ingest reviewed} (\texttt{education\_ingest.py}; the factorial question for this
        domain is still open and deliberately untouched).
        \begin{itemize}
          \item \textbf{Bug: the essay-id lookup ignored its own priority list} and scanned the file's
                column order instead, so it keyed on \texttt{essay\_id} --- 661 of whose 15{,}593 values
                are Excel-mangled to scientific notation (\texttt{5.88194E+12}), one of them covering two
                different essays, which the per-essay dedup then silently merged. It now prefers
                \texttt{essay\_id\_comp}: 15{,}594 distinct essays.
          \item \textbf{PERSUADE's real demographics are now kept on the record} (sex, race/ethnicity,
                school grade; ELL, economic and disability status in \texttt{extra}) and still never
                rendered, as in credit. They are per-essay consistent across the discourse rows. This is
                the only domain where all three injected axes have a real counterpart.
          \item \textbf{Both loaders now return a filter report} (as \texttt{credit\_clean} /
                \texttt{bios\_clean} do), logged into the manifest. It makes two things visible: the
                PERSUADE length filters are nearly inert (0 essays below 300 characters, 58 above 6{,}000;
                the real cut is the 9{,}132 middle scores), and ASAP's ``tercile'' split is not a third ---
                the scores are coarse integers, so only 574 of 12{,}976 essays fall in the middle and
                6{,}424/4{,}782 are labelled strong/weak.
          \item ASAP no longer reports its essay set as a \texttt{grade\_level}: the field means a real
                school grade or nothing, and the set stays in \texttt{prompt\_id}/\texttt{extra}.
          \item \textbf{Education is a factorial too: sex $\times$ ethnicity $\times$ economic status},
                so all three domains now share one design shape. Pole A is the hypothesised penalised
                level on every axis (female, Black, low income).
                \begin{itemize}
                  \item \textbf{Why economic status and not stage.} A stated stage contradicts the
                        essay's register and has no signed direction (norm-referenced grading and status
                        credibility pull opposite ways), so the corner cell could not be called ``most
                        penalised''. An income level contradicts nothing: 5 of 6{,}397 essays mention
                        their own household money at all, so the factorial needs no plausibility filter
                        and keeps the whole corpus (6{,}392) at its 40\% strong rate, instead of the
                        stage design's 3{,}648 at 25\%. The stage axis keeps its own manifest and
                        population; every education result must say which of the two it used.
                  \item \textbf{Sex and ethnicity share one carrier} --- both proxies are the first name.
                        So the proxy encoding is the Haim sex $\times$ ethnicity name grid drawn at a
                        single pool index (\texttt{ProxyNames.draw\_grid}): index-matched, so a sex or
                        ethnicity swap moves one step in the grid rather than resampling two independent
                        names and charging the difference to the axis.
                  \item \textbf{No sex noun in the explicit clause} (``The student is Black, female, and
                        from a low-income household''): ``girl'' vs ``young woman'' would make the sex
                        wording covary with the writer's age --- the error the hiring career-gap proxy
                        made. Three independent slots instead.
                  \item \textbf{Economic proxy $=$ the school's free/reduced-price-lunch share} (``most''
                        vs ``few'' students qualify), the standard US school-poverty measure, phrased
                        without a negation. It is school-level, so the proxy measures school poverty
                        where the explicit clause states household income; labels
                        \texttt{high\_poverty\_school}/\texttt{low\_poverty\_school} say so.
                  \item Worst-case clause length delta is \textbf{5 characters} (proxy ethnicity; the
                        composed corner is 4), so every pair passes the Tier-1 gate at the default
                        bounds --- no relaxed intersection threshold, unlike credit and hiring, and 0 of
                        1{,}200 blocks dropped in a 300-essay run.
                  \item \textbf{Sign bug found by the port.} Education's \texttt{ethnicity} was declared
                        \texttt{pole="b"} in the decision-response frame, because the old single-axis
                        marker had white as pole A. Under the factorial that pointed the decision arm at
                        the \emph{reference} pole, i.e.\ it would have measured the discriminatory verdict
                        against a white student and reported it as the protected-pole result. Fixed and
                        pinned by a regression test.
                  \item Real writer attributes (sex, race/ethnicity, grade, ELL, economic and disability
                        status) go into \texttt{cells.jsonl} as \emph{covariates}, not factors: they play
                        no part in detection or debiasing (the matched pair controls them by
                        construction) and belong to the validity section. Their one load-bearing use is
                        the RQ4 check --- after nulling, is the quality signal degraded together with the
                        real-group signal? --- since in this corpus quality is badly entangled with
                        group (ELL 1.78 vs 2.64; economically disadvantaged 2.06 vs 3.15).
                \end{itemize}
          \item \textbf{The stage axis (\texttt{grade\_level}) was widened to 6th grade vs doctoral
                candidate}, and the intermediate rungs (12th grade, final-year bachelor's, master's,
                doctorate) run as a separate \emph{monotonicity sweep}: every rung is contrasted against
                the same 6th-grade reference clause, so one battery run over
                \texttt{stage\_grade12,\dots,stage\_doctorate} traces a curve. A smooth monotone curve
                reads as norm-referenced grading; a jump at a status boundary reads as a prior about the
                writer. The top rung duplicates \texttt{grade\_level} as a consistency check.
                \begin{itemize}
                  \item \textbf{Plausibility is mostly a prompt property, not a text property.} The six
                        PERSUADE prompts that ask pupils about their own schooling draw first-person
                        pupil talk in 46--67\% of essays; the source-text prompts draw it in 0--10\%
                        (``The Face on Mars'' 0.0\%). So the filter selects the nine stage-neutral
                        prompts first, then applies cue rules (own school life, a named stage, own grade,
                        pupil voice, school routine, ``Dear Principal'' letters) as a safety net:
                        6{,}404 essays $\to$ \textbf{3{,}648}, of which 2{,}459 go to prompt selection.
                        ``college''/``university'' mentions are deliberately not a cue.
                  \item \textbf{The low pole is the corpus's own floor} (PERSUADE has no grade below 6),
                        so some essays genuinely are that young. What no filter can fix is \emph{register}:
                        the essays read like grades 6--10 whatever the marker claims.
                  \item \textbf{Grading is norm-referenced by nature}, and widening the range maximises
                        that: told the writer is in 6th grade, a competent grader should score a fixed
                        essay \emph{higher}, which points opposite to the bias hypothesis. Read the stage
                        gap together with the quality interaction --- norm-referencing predicts a gap that
                        grows with essay quality, a flat offset is a prior about the person. Whether to
                        pin the rubric in the assessment prompt (``against a fixed standard, regardless of
                        who wrote it'') or to run both framings as a factor is still open.
                  \item The proxy clauses are stage-only (``is enrolled in \dots''), never an achievement
                        or duty: ``teaches an undergraduate section'' would be an accomplishment cue, not
                        an age cue --- the mistake the hiring family proxy made with a career gap. All
                        rungs sit within 2 characters of the reference clause, so the Tier-1 gate passes
                        with no relaxation (0.00 discard on every cell).
                  \item \textbf{To watch:} the neutral prompts are graded harder --- only 913 of the 3{,}648
                        are ``strong'' (25\%, vs 40\% before), so the \crossinf{} quality balance is
                        thinner; and they cover grades 6--10 only, with 10 the largest (1{,}708) and 6 the
                        smallest (364), which is the pair for the real-grade confound check.
                \end{itemize}
          \item \textbf{The grader now sees the assignment.} The header used to be
                ``Student essay submission.'' plus the body, so the RM graded an answer to a question it
                could not see --- and the stage-neutral prompts are mostly \emph{text-dependent}: the
                essay argues about a source article, and 977 of the 3{,}647 kept essays (27\%) say ``the
                article''/``the author'' outright (``Exploring Venus'' 74\%). The task text is constant
                per prompt (199--598 characters against a median essay of 2{,}135) and identical across an
                A/B pair, so it is rendered into the header exactly as the hiring header states the target
                role, and omitted where the corpus has none (ASAP). This does not touch the bias contrast;
                it repairs the \emph{quality} arm, since ``does it address the prompt'' is a rubric
                dimension the model could not apply --- and education is the domain \crossinf{} was
                supposed to be interpretable in.
          \item \textbf{Paragraph structure is no longer flattened.} 15{,}004 of the 15{,}594 essays are
                paragraphed (median 4 blank-line breaks) and the cleaner collapsed all of it into one
                line. Blank lines now survive; everything else still collapses to a single space, so a
                hard-wrapped corpus would not gain false breaks (PERSUADE has no lone newlines). Same
                reasoning as above: identical on both sides, so no contrast was biased, but every essay
                looked structureless. Side effects: bodies are a few characters longer, so 65 rather than
                58 exceed the length cap (6{,}397 essays kept, 3{,}647 after the stage rules), and A2's
                insertion points no longer include sentences that end a paragraph.
          \item Cosmetic but real: the \texttt{csv} field-size limit is no longer raised as an import side
                effect; the ASAP pass no longer writes floats into the reader's row dict; the essay body is
                cleaned once per essay rather than once per discourse row. The @-redaction regex also fires
                on the 5 PERSUADE essays containing a bare handle (``@TimmyTurner''), which is wanted for a
                name and coarse otherwise --- kept, and now counted in the report.
        \end{itemize}
\end{itemize}

\todo{Numbers above are literal/preliminary --- wire into \texttt{shared/numbers.tex} via
\texttt{export\_paper\_numbers.py} once stable; promote settled results into progress\_*.tex, then sections/.}

\bibliographystyle{plainnat}
\bibliography{shared/references}
\end{document}
